\documentclass[letterpaper,11pt]{article}
\usepackage{times}
\usepackage[margin=1in]{geometry}
\usepackage[singlespacing]{setspace}
\usepackage[hidelinks]{hyperref}
\usepackage{graphicx}
\usepackage{booktabs}
\usepackage{array}

\begin{document}

% ===================== TITLE PAGE =====================
\begin{titlepage}
\begin{center}
\large
BOSTON UNIVERSITY\\[0.5em]
DEPARTMENT OF ELECTRICAL AND COMPUTER ENGINEERING\\[2em]
PhD Prospectus\\[3em]

{\bfseries MULTI-LAYER SECURITY FOR MOBILE USERS:
MEASUREMENT AND MITIGATION FROM APPLICATIONS TO CONVERSATION}\\[3em]
By\\[2em]

OLAWALE AMOS AKANJI\\
B.Sc.\ Cyber Security, Air Force Institute of Technology, Nigeria, 2023\\
M.S., Boston University, 2026\\[3em]

Advisor: Prof.\ Gianluca Stringhini\\
Co-Advisor: Prof.\ Manuel Egele
\end{center}

\vfill

\begin{center}\bfseries Abstract\end{center}
\noindent
Mobile users face security and privacy threats at every layer of the
ecosystem they depend on: the applications they install, the platform
mechanisms that govern those applications, and the conversations they
hold on their devices. This dissertation develops defenses for users at
each of these layers. At the application layer, \emph{The Cost of
Convenience} (ACM ASIA CCS 2026) presents the first cross-country
measurement of predatory loan-app compliance with national regulations
and Google's Financial Services Policy, analyzing 434 apps across five
countries; following our disclosures, Google removed 93 apps
representing over 300 million installs. Investigating how such apps
retain access to sensitive data despite policy restrictions led to the
platform layer: \emph{Silent Consent, Persistent Risk} shows, through a
longitudinal analysis of 19.3 million APKs, that Android's permission
groups silently auto-grant new permissions across updates (17\% of
multi-version apps expand silently, and malware does so at
significantly higher rates) and that normal-level custom permissions
expose sensitive data to unrelated apps without any prompt; a
lightweight prototype restores update-time transparency without OS
modification. Yet even compliant apps on a hardened platform cannot
protect users who are manipulated directly. The third layer is the
conversation. \emph{SIRENBERT} detects \emph{conversational scams},
multi-turn exchanges in which a suspect under a false identity,
situation, or purpose induces the target to part with money, by
modeling the exchange message by message on the user's device: a
compact encoder labels each message with one of 14 observable
behavioral triggers, and a recurrent model over the trigger sequence
tracks the conversation's lifecycle and raises a warning before the
financial demand. On 22,750 conversations from 5 real and 5 synthetic
sources, labelled message by message under a sequential constraint,
the trigger stream carries information the demand alone does not:
real scam conversations show a median of 2 public-guidance indicators
before the first demand, benign conversations with a demand show none
in 90\% of cases. The proposed work extends the conversation layer
beyond text: scammers send images, screenshots, voice notes, and video
whose function is credibility, and the dissertation will measure these
\emph{credibility artifacts} at scale, test whether their claims agree
with the conversation's, and fold them into the on-device trajectory
model. Together the three layers defend mobile users wherever
exploitation reaches them.
\end{titlepage}

% ===================== BODY =====================

\section{Introduction}
\label{sec:intro}

Mobile devices mediate nearly every aspect of their users' lives, and
for billions of people, particularly in emerging markets, they are the
sole gateway to financial services~\cite{suri2016mpesa,findex2021}.
This centrality exposes users to security and privacy threats at
multiple layers of the mobile ecosystem: the \emph{applications} they
install, the \emph{platform mechanisms} that are supposed to govern
what those applications can do, and the \emph{conversations} they hold
through their devices. My research protects users at each of these
layers. Rather than one system, this dissertation comprises three
studies, each identifying how users are harmed at one layer and
delivering a concrete, deployed or deployable defense.

My work began at the application layer. In emerging markets, digital
lending applications (loan apps) have become a primary channel for
microcredit, but many demand excessive permissions and weaponize the
harvested data, contact lists, SMS, photos, for coercive debt
recovery: harassment, blackmail, and public shaming of borrowers and
their contacts. In \emph{The Cost of Convenience: Identifying,
Analyzing, and Mitigating Predatory Loan Applications on
Android}~\cite{akanji2026cost}, we conducted the first cross-country
measurement of loan-app compliance with national regulations and
Google's Financial Services Policy, and our disclosures led Google to
remove 93 apps with over 300 million cumulative installs
(\S\ref{sec:loanapps}).

That analysis raised a deeper question. While studying how loan-app
developers retain access to sensitive data despite policy
restrictions, we observed that circumvention does not always require
malice or obfuscation: certain Android mechanisms enable it by
default. Once a user approves one permission from a group, updates can
silently add sibling permissions without any prompt, and developer-defined
custom permissions at the default protection level are auto-granted at
install with no user visibility at all. This moved my focus to the
platform layer. In \emph{Silent Consent, Persistent Risk: Android
Permission Groups and Custom Permissions}~\cite{akanji2026silent}, we
measured both mechanisms at ecosystem scale, 19.3 million APKs,
validated their exploitability on Android 16, showed that malware
preferentially engages them, and built a transparency prototype that
restores update-time consent visibility using only public APIs
(\S\ref{sec:silentconsent}).

Applications and the platform, however, are not the only sources of
risk: users can be exploited directly, through the conversations their
devices host. In 2025 the FBI's Internet Crime Complaint Center
received 1,008,597 complaints totaling \$20.9 billion, and the
categories that operate through sustained conversation account for
most of it: investment fraud alone \$8.6 billion, confidence and
romance schemes a further \$929 million~\cite{ic32025}. No permission
policy or platform hardening can protect a user who willingly sends
money to someone they trust. My third study, \emph{SIRENBERT}, defends
this layer for \emph{conversational scams} as a class, romance,
investment, advance-fee, employment and related persuasion-based fraud
alike, with a privacy-preserving, on-device detector that models the
conversation's behavioral trajectory message by message
(\S\ref{sec:sirenbert}). The proposed work extends that defense to the
images, screenshots, voice notes, and video that scammers send to
manufacture credibility (\S\ref{sec:proposed}).

\textbf{Thesis statement.} \emph{Mobile users are exploited at three
distinct layers, the applications they install, the platform mechanisms
that govern those applications, and the conversations their devices
host, and a defense at one layer does not protect against the others.
Protecting users therefore requires measuring the harm and deploying a
defense at each layer, and this dissertation does both.}

\section{Background and State of the Art}
\label{sec:background}

\subsection{Predatory Digital Lending}
Loan apps promise fast, collateral-free microcredit to populations
underserved by formal banking~\cite{suri2016mpesa}. Many, however,
obscure predatory terms and adopt digitized coercive debt-collection
tactics enabled by data harvested at install time. Governments in
affected countries have introduced registration and privacy
constraints for digital lenders, and Google's Financial Services
Policy (FSP) prohibits personal-loan apps from requesting high-risk
permissions such as \texttt{READ\_CONTACTS} and
\texttt{READ\_SMS}~\cite{googlefsp}; yet abuse reports, and user
suicides linked to loan-app harassment, persist. Prior academic work
is primarily qualitative, interviews and case studies documenting
distress and perceived data misuse; before our work there was no
large-scale technical measurement of whether regulator-approved apps
actually comply with either national rules or platform policy.

\subsection{The Android Permission Model and Its Legacy Mechanisms}
Android's permission system has been studied extensively for
over-privilege and user
comprehension~\cite{felt2011demystified,felt2012attention,wijesekera2015remystified}.
Android 6.0's runtime permissions were designed to restore informed
consent, but two mechanisms inherited from Android 1.0 remain active
in Android 16. \emph{Permission groups} reduce prompt fatigue by
auto-granting new permissions within an already-approved group during
updates, with no dialog and no Settings indication. \emph{Normal-level
custom permissions}, the default when developers omit a protection
level, are auto-granted at install to any requesting app regardless of
developer identity. Prior work documented confused-deputy attacks and
isolated exposures~\cite{tuncay2018custom,gamba2020preinstalled}, and
Calciati et al.~\cite{calciati2020automatically} measured group
auto-grant on 2.9M app versions, but the mechanisms' current
prevalence, their association with malice, and their exploitability on
modern Android remained open questions.

\subsection{Conversational Scams}
A \emph{conversational scam} is a multi-turn exchange in which a
suspect uses a deceptive identity, situation, or purpose to induce the
target to part with financial resources, through a request or
instruction whose fraudulent nature is supported by earlier evidence in
the conversation. Romance fraud, pig butchering, advance-fee,
employment, and technical-support fraud are its families; Sun et
al.\ organise 177,989 victim reports into 11,573 such instances across
20 categories~\cite{sun2026prescam}. Studies of individual families
describe the same progression in different vocabularies: contact,
trust building, pre-exploitation setup, and a financial
demand~\cite{whitty2013scammers,oak2025anna,acharya2024explorative,spokoyny2025victim}.
The setup phase is built to be indistinguishable from legitimate
conversation, which defeats existing detection: keyword and rule
detectors fire only on money language, after the grooming, and misfire
on ordinary money talk; learned classifiers judge completed
transcripts, after the harm; prompted language models are accurate on
whole conversations but too costly and privacy-invasive to run on every
private message, and reported error rates on romance and wrong-number
exchanges are substantial~\cite{lu2026readthis,topcuoglu2026genai}.
Before this work no labeled dataset followed these conversations
message by message, the granularity at which a deployable detector
must operate, and no detector treated the class as a whole.

\section{Completed Work I: The Cost of Convenience}
\label{sec:loanapps}

\emph{The Cost of Convenience: Identifying, Analyzing, and Mitigating
Predatory Loan Applications on Android}~\cite{akanji2026cost},
published at ACM ASIA CCS 2026, protects users at the application
layer through measurement, disclosure, and policy recommendations.

\textbf{Approach.} We collected 434 loan apps from regulator registries
and app markets in Indonesia, Kenya, Nigeria, Pakistan, and the
Philippines, translated regulatory and platform policy prose into
testable permission checks with an LLM-assisted
\emph{policy-to-permission mapping}, and combined it with static
analysis of app code and dynamic analysis of what data leaves the
device and when.

\textbf{Key findings.} Non-compliance is pervasive among
\emph{approved} apps: 141 violate national regulatory policy and 147
violate Google's FSP, with combined downloads above 300 million; 37 apps
transmit contacts, SMS, location, or media at launch, before any signup,
eliminating even the possibility of informed consent. National
regulations and platform policy are systematically misaligned, so an
app can satisfy one authority while violating the other.

\textbf{Impact.} Following our coordinated disclosure, Google removed
93 flagged apps from Google Play. The methodology is directly usable
by regulators and platforms as a proactive compliance-monitoring
pipeline.

\section{Completed Work II: Silent Consent, Persistent Risk}
\label{sec:silentconsent}

Studying how developers circumvent permission policies in
\S\ref{sec:loanapps} revealed that the platform itself provides
circumvention paths by default. \emph{Silent Consent, Persistent Risk:
Android Permission Groups and Custom
Permissions}~\cite{akanji2026silent} measures these paths at ecosystem
scale and protects users by restoring the visibility the platform
withholds.

\textbf{Research questions and approach.} How far do permission groups
enable silent capability expansion across updates; is that expansion an
ecosystem-wide consent failure or disproportionately a property of
malware; how prevalent are normal-level custom permissions that expose
data across developers, and is exploitation confirmed on current
Android; and can an interface-level intervention restore update-time
visibility without OS modification? We analyzed 19,324,760 APKs from
5,965,583 apps in AndroZoo~\cite{allix2016androzoo}, tracking
permission evolution across versions, operationalized malice through
VirusTotal detections under a threshold sensitivity
sweep~\cite{zhu2020vt}, validated exploitability with per-pair probe
apps on Android 16, and attributed the exposed surface to app-core code
versus bundled libraries.

\textbf{Key findings.} Among 2,244,575 multi-version apps, 381,026
(17\%) silently gain permissions within already-granted groups, and
the auto-grant fires across every dangerous group on Android 16 with
no runtime dialog. Silent expansion is preferentially engaged by
malware: odds ratio 1.35 ($p<0.001$), rising to 2.06 in the most
permission-heavy quartile, and significant at every tested detection
threshold. We identified 307 cross-developer normal-custom-permission
pairs re-delegating contacts, SMS, location, authentication
credentials, identity, and medical records to unrelated apps without
any prompt, confirmed active exploitability with instrumented probe
apps, and attributed most of the exposed surface to bundled
third-party libraries rather than app-core code.

\textbf{Mitigation.} A prototype built on public Android APIs restores
update-time transparency: in a 96-day deployment it surfaced 23 silent
expansion events across 13 apps, including apps with over one billion
installs, at about one notification every four days.

\section{Current Work III: SIRENBERT}
\label{sec:sirenbert}

\emph{SIRENBERT: On-Device Detection of Conversational Scams by
Modeling Messages in Conversations} defends the conversation layer.
The system is built and deployed in an Android messaging client; the
remaining work is the re-annotation, retraining, and evaluation
described in \S\ref{sec:tasks}.

\textbf{Definition and characteristics.} We adopt the definition in
\S\ref{sec:background} and operationalise its three aspects. A
multi-turn exchange: every conversation in our data has at least 7
messages; real scam conversations have a median of 66 and reach the
first financial demand at median message 17, and the demand opens the
conversation in 1 case out of 4,298. Deception: observed through
misrepresentation signals in the conversation, never through a
judgement of what is true. A financial demand: a message whose
compliance would cost the target money or economic value. Across the
real scam conversations no persuasion tactic is universal, emotional
appeal appears in 84\%, romance in 61\%, an investment offer in 23\%,
which is why the class, and the detector, are defined over the
sequence and not over any one tactic.

\textbf{Lifecycle and triggers.} The progression documented by the
family studies is formalised as 5 lifecycle states over 14 observable
\emph{behavioral triggers}: 2 contact behaviors, 3 relationship
behaviors, 8 pre-exploitation behaviors (credibility positioning, gain
offers, verification avoidance, platform migration, love bombing,
crisis disclosure, emotional manipulation, urgency pressure), and the
financial demand. Each trigger names what a message does to the
conversation, and a trigger is added to the set only if it changes the
next lifecycle transition. A demand carries a sub-type, exploitative or
benign, and a structured basis citing the earlier messages that make
it exploitative; a conversation is labelled a scam only when an
exploitative demand cites earlier pre-exploitation behavior, derived by
one versioned rule from the message labels.

\textbf{Dataset.} No public dataset follows scam conversations message by
message, so we assembled a proxy dataset from sources that each capture
a different part of the phenomenon (Table~\ref{tab:data}). LoveFraud02
contributes 83 romance-fraud conversations recorded with genuine scammers
by a researcher, 30 to 3,405 messages each. Reddit Scambait contributes
1,142 conversations recovered by OCR from screenshots posted to
r/scambait and verified against the images; the scammer side is
authentic and the pretexts vary widely, which makes it our main source
of family diversity, and the scambaiter side is a distribution shift we
evaluate explicitly. Topical-Chat and XDailyDialog supply benign
dialogue, the latter chosen because its warm interpersonal tone is close
to the surface of romance scams. PersuasionForGood supplies 1,017
charity solicitations in which one person persuades another to give
money: sustained, emotionally framed, and directed at a genuine demand,
the closest available match to what a detector must not flag. Two
generated corpora, a pig-butchering set built to the family's full
extraction cycle and ScamForge conversations grounded on 3,288 BBB
romance-scam reports, cover behavior combinations rare in the real
sources and are used for stress testing only, never in the primary test
set. Every conversation is labelled message by message by a language
model annotator under a sequential constraint (it sees only the prefix a
deployed detector would see), with a second annotator over the same
conversations to measure annotator dependence. Real scam conversations
are 48\% romance, 9\% romance-investment (pig butchering), 11\%
shopping and delivery, 7\% friend-or-family emergency, 6\% employment,
5\% advance fee; 1,916 real benign conversations contain a financial
demand. Sextortion is outside the scope and held separately.

\begin{table}[h]\centering\footnotesize\setlength{\tabcolsep}{4pt}\begin{tabular}{@{}lrrr@{}}\toprule
Source & Type & Convs. & Annotated msgs. \\\midrule
LoveFraud02 (victim transcripts, romance)~\cite{faber2024lovefraud02} & real scam & 83 & 26,423 \\
Reddit Scambait (OCR of screenshot posts) & real scam & 1,142 & 65,685 \\
Topical-Chat~\cite{gopalakrishnan2019topical} & real benign & 8,628 & 97,204 \\
XDailyDialog~\cite{li2017dailydialog} & real benign & 3,100 & 19,981 \\
PersuasionForGood~\cite{wang2019persuasion} & real benign & 1,017 & 10,600 \\
Generated pig butchering (scam / benign) & synthetic & 5,683 & 237,296 \\
Generated ScamForge (scam / benign / hard negative) & synthetic & 3,097 & 152,371 \\
\midrule \textbf{Total} & & \textbf{22,750} & \textbf{609,560} \\\bottomrule\end{tabular}
\caption{The SIRENBERT dataset after the sextortion exclusion. Annotated messages are the suspect side, the side a detector classifies.}\label{tab:data}\end{table}

\textbf{Lifecycle states and triggers.} Table~\ref{tab:triggers} gives
the vocabulary. The state after a message is a deterministic function of
the previous state and the trigger: contact and small talk reach at most
\textsc{initial\_contact}, relationship triggers
\textsc{relationship\_building}, pre-exploitation triggers
\textsc{pre\_exploitation}, and a financial demand reaches
\textsc{exploitation} only from \textsc{pre\_exploitation}; states
never move backward. The state summarizes what has been observed; the
scam verdict is separate: a conversation is \textsc{scam} when an
exploitative demand cites earlier pre-exploitation behavior,
\textsc{suspicious} when such behavior occurs with no qualifying
demand, and \textsc{non\_scam} otherwise, derived by one versioned rule
from the message labels and never from the annotator's own opinion.

\begin{table}[h]\centering\footnotesize\setlength{\tabcolsep}{4pt}\begin{tabular}{@{}llp{3.6in}@{}}\toprule
State admitted & Trigger & Applies when \\\midrule
Initial contact & WrongNumber\_Intro & contact opens under a wrong-number pretext \\
 & Introduction\_Opener & a self-introduction to a stranger \\\midrule
Relationship & Trust\_Building & personal history shared or invited to prompt disclosure \\
building & Affectionate\_Language & warmth within ordinary courtship \\
 & SmallTalk\_Maintenance & routine conversation that sustains contact \\\midrule
Pre-exploitation & Credential\_Positioning & standing, wealth, authority or expertise asserted for credibility \\
 & Gain\_Offer & an investment, job, prize or other gain held out to the target \\
 & Verification\_Avoidance & identity checks deflected, or suspect-controlled proof presented \\
 & Platform\_Migration & moving the conversation to another channel \\
 & Love\_Bombing & romantic love or commitment declared to the target \\
 & Crisis\_Disclosure & an acute emergency presented \\
 & Emotional\_Manipulation & guilt, obligation or affection used as leverage \\
 & Urgency\_Pressure & deadlines or scarcity imposed \\\midrule
Exploitation & Financial\_Demand & a request whose compliance costs the target money or economic value; carries a sub-type and a cited basis \\
\bottomrule\end{tabular}
\caption{The 14 behavioral triggers. Each names what a message does to the conversation, applied to the text alone; none asserts intent.}\label{tab:triggers}\end{table}

\textbf{Architecture.} Detection is two stages on the device.
\emph{Stage 1}, inside the messaging application, reads message $m_t$
with the $k$ preceding messages, each prefixed by its sender, and emits
a probability over the 14 triggers and, for a demand, its sub-type; we
compared DistilBERT and ModernBERT at 512 and 8,192 tokens across
$k \in \{0,\ldots,15\}$, and a rank-64 LoRA adapter over a frozen
Gemma-4-E4B, the class of model phones already ship, which reaches
$0.986 \pm 0.005$ conversation SCAM-F$_1$ where prompting the same model
fails at 0.06 macro-F$_1$. \emph{Stage 2}, a unidirectional GRU over the
$[T \times 14]$ trigger sequence, re-scores the conversation after every
message with a terminal head for the scam decision and an auxiliary
head for whether setup has begun; it is unidirectional by requirement,
since a detector that runs as messages arrive cannot condition on the
future. Only trigger vectors are buffered; message text is never stored
or transmitted. The pipeline runs in Element X for Android through
ONNX Runtime Mobile at 435\,ms per message; export is lossless, dynamic
INT8 quantization is the entire on-device accuracy gap (recall 0.804 to
0.686), so FP16 is the deployment artifact.

\textbf{Evaluation protocol.} Model selection uses nested 5-fold
cross-validation over pooled training data with the test partition
untouched~\cite{cawley2010overfitting,varma2006bias}; every reported
number is a mean over 3 seeds with bootstrap 95\% intervals; results
are reported per source as well as pooled, and entire sources and
families are held out to test transfer. Detection lead time relative to
the first validated demand is the deployment-critical metric, and every
deployment claim is backed by bulk replay of the held-out partition
through the Android client on real hardware.

\textbf{Results to date.} On the evidence side, real scam conversations
show a median of 2 public-guidance indicators before the first demand
(51.8\% show 2 or more), while benign conversations with a demand show
a median of 0 and reach the demand with none in 90.1\% of cases;
relationship-building indicators rise steeply with conversation length
(claims of affection 11\% to 82\% across length quartiles) while the
investment offer falls. On the model side, the nested-cross-validation
selection reaches 0.863 $\pm$ 0.005 conversation SCAM-F$_1$ over 3
seeds on a 1,297-conversation held-out partition under the labels it was
trained on; under the revised rule-derived labels the same models score
0.70, because the partition gains 22 positives they never saw, which is
why retraining under the rule is the first remaining task.

\subsection{Remaining Tasks for SIRENBERT}
\label{sec:tasks}

\textbf{RT1 --- Re-annotation under the revised guide and retraining.}
Version 7 of the annotation guide records structured evidence
citations, drops the intent field, and quarantines invalid output;
after a 200-conversation pilot with human review, the full corpus is
re-annotated and the encoders and recurrent model retrained under the
rule-derived labels, with nested cross-validation, 3 seeds, and
bootstrap intervals as before.

\textbf{RT2 --- Generalisation and robustness.} Leave-one-family-out
evaluation (train on romance and investment, test on advance fee and
employment, and the converse) tests whether the trigger stream
transfers across families; the adversarial suite (distractor padding,
paraphrase, slow-play stretching of the lifecycle) quantifies robustness
bounds. A facet ablation on two triggers decides whether finer
distinctions add detection value.

\textbf{RT3 --- Cross-application state and user-facing warnings.}
Platform migration is a scripted step, so the exploitation often
completes in a different app from the grooming; the on-device design
will carry trajectory state across applications under explicit consent,
and an IRB-approved study will evaluate warning timing and framing with
users who are already emotionally invested.

\section{Proposed Work: Credibility Artifacts in Conversational Scams}
\label{sec:proposed}

SIRENBERT reads text, but scammers also send a photo in uniform, a
passport, a screenshot of a trading balance, a voice note, a video call
with a borrowed face. The behavior these serve, presenting
suspect-controlled evidence as proof or deflecting verification, appears
in 63\% of real scam conversations, and the artifact itself is
invisible to a text detector. The proposed work studies these
\emph{credibility artifacts} as part of the conversation layer.

\textbf{RQ1 --- Measurement.} What artifacts do scammers send, at which
lifecycle stage, and how much of it is reused or synthetic? The Reddit
scambaiting corpus consists of chat screenshots that embed the images
scammers sent; a new crawl retaining full-resolution originals and post
metadata, followed by extraction of the embedded images, yields the
first large-scale corpus of scam artifacts with their conversational
context. Reuse is measured by perceptual hashing against known
scam-photo collections and across our own corpus; synthetic share by
generation detectors, reported with their known decay over time.

\textbf{RQ2 --- Cross-modal consistency.} Does the artifact's claim
agree with the conversation's? A text that says oil-rig engineer beside
a stock photograph of an army officer is a contradiction no single
modality can see. We will extract the claims an artifact makes
(identity, occupation, location, wealth, proof of returns) and test
them against the triggers already recorded for the conversation,
turning an artifact into an event in the same trajectory.

\textbf{RQ3 --- On-device fusion.} An artifact event becomes a trigger
with its own evidence, consumed by the same recurrent model. The
privacy constraint is argued per artifact type: perceptual-hash lookups
can be made private, uploading a photograph cannot, and voice notes are
transcribed on the device and enter the text pipeline unchanged. We
will measure what each artifact type adds to detection lead time and
false-alarm rate over the text-only system. Images, screenshots, and
voice notes are in scope; live video calls are treated as a threat and
evaluated only where public material exists.

\subsection{Anticipated Contributions}
\label{sec:contrib}

This dissertation will contribute, at the application layer, the first
cross-country compliance-measurement methodology for financial apps,
with demonstrated real-world impact (93 removals, 300M+ installs); at
the platform layer, the largest longitudinal measurement of consent
erosion in Android permissions, evidence that malware preferentially
exploits it, and a practical transparency countermeasure requiring no
OS modification; and at the conversation layer, a definition and lifecycle
formalization of conversational scams as a class, the first
message-level labeled dataset for it, an on-device early-warning
detector evaluated across scam families, and the first large-scale
measurement and detection of the credibility artifacts scammers send. Collectively, the work
advances mobile security from post-hoc, transcript-level, single-point
responses toward proactive protection of users at every layer where
exploitation reaches them.

\section{Research Plan and Timeline}
\label{sec:timeline}

\begin{table}[h]
\centering
\small
\begin{tabular}{p{1.4in}p{4.6in}}
\toprule
\textbf{Period} & \textbf{Milestones} \\
\midrule
Fall 2026 & Prospectus defense. RT1: v7 pilot, full re-annotation, retraining under the
rule-derived labels; submit SIRENBERT. New Reddit crawl with originals and metadata. \\
Spring 2027 & RT2: leave-one-family-out and adversarial suite. RQ1: artifact
extraction and measurement (reuse, synthetic share, lifecycle stage). \\
Summer 2027 & RQ2: cross-modal consistency; RT3 study design and IRB approval. \\
Fall 2027 & RQ3: on-device fusion and evaluation; submit the artifacts paper; RT3 user study. \\
Spring 2028 & Complete remaining evaluations; dissertation writing and defense. \\
\bottomrule
\end{tabular}
\caption{Anticipated timetable for completion.}
\end{table}

\section{Summary}

Mobile users are exploited through the apps they install, the platform
mechanisms meant to govern those apps, and the conversations they
trust. My completed work has protected users at the first two layers,
removing predatory apps from the ecosystem and restoring visibility
into silent consent erosion, each with demonstrated real-world impact.
The third layer is defended by SIRENBERT, an on-device early-warning
system for conversational scams that reaches users before the financial
demand; the proposed work extends it to the images, screenshots, and
voice that scammers use to manufacture credibility.

% ===================== BIBLIOGRAPHY =====================
{\footnotesize
\bibliographystyle{plain}
\bibliography{prospectus_revised_draft}}

% ===================== CV =====================
% Attach up-to-date CV here (not counted toward the 10-page limit).

\end{document}
